% !TeX spellcheck = it_IT

\documentclass[handout]{beamer}

\usepackage[utf8]{inputenc}
\usepackage[italian]{babel}
\usepackage{multicol}
\usepackage{tabularx}
\usetheme{metropolis}

\usepackage{mathtools}
\DeclarePairedDelimiter{\ceil}{\lceil}{\rceil}


\title{Informatica B - Esercitazione 0}
\subtitle{Codifica binaria - ulteriori esercizi}

\author{Stefano Cereda\\
		stefano1.cereda@mail.polimi.it
	}
\date{26/09/2017}
\institute[PoliMi]{\vspace{0.5cm}\centering Politecnico di Milano \\ \vspace{0.2cm}
	\includegraphics[width=0.2\linewidth]{./logopolimi}}
%\logo{\includegraphics[width=15mm]{./logopolimi}}

%\setbeamercovered{invisible}

\makeindex

\begin{document}
\begin{frame}
	\maketitle
\end{frame}

\begin{frame}{Codifica posizionale}
La codifica che utilizziamo tutti i giorni è detta posizionale in quanto ogni cifra assume un significato diverso a seconda della posizione occupata:

$23_{10} = 2*10^1 + 3*10^0 = 20 + 3$

Questo meccanismo può essere utilizzato per trasformare in base dieci un numero espresso in qualsiasi altra base (purchè utilizzi una codifica posizionale semplice):
$110_2 = 1*2^2+1*2^1+0*2^0 = 4 + 2 +0 = 6_{10}$
\end{frame}

\begin{frame}{Esempi conversione da base 2 a base 10}
Ricordando le potenze di 2 (1,2,4,8,16,32,64,128,...) la conversione diventa molto più facile. Basta sommare le potenze in corrispondenza degli uni:

$01101101_2 = 1+4+8+32+64 = 109_{10}$

$10010010_2 = 2+16+128 = 146_{10}$
\end{frame}

\begin{frame}{Basi 8 e 16}
Altre basi molto utilizzate in ambito informatico sono la base 8 e la base 16. Questo perchè si può facilmente convertire in base 8 un numero in base 2 considerando blocchi di tre cifre:
$110011_2 \rightarrow 110 \ 011 \rightarrow 6 \ 3 \rightarrow 63_8$
$011110_2 \rightarrow 011 \ 110 \rightarrow 3 \ 6 \rightarrow 36_8$\\
Ovvero, la tripletta $011$ si traduce sempre in un $3$.

Allo stesso modo, possiamo convertire da base 2 a base 16 considerando quattro bit:
$01101101_2 \rightarrow 0110 \ 1101 \rightarrow 6 \ 13_{10} \rightarrow 6 \ D \rightarrow 6D_{16}$

Le cifre della base 16 sono: 0,1,2,3,4,5,6,7,8,9,A,B,C,D,E,F
\end{frame}

\begin{frame}{Conversione da base 10 a base 2}
Per convertire un numero da base 10 a base 2, dobbiamo dividere il numero fino ad arrivare a 0. Nel farlo teniamo traccia dei resti: letti al contrario ci daranno la codifica in base 2.

\begin{tabular}{c|c}
		12 & \\
\end{tabular}
$\rightarrow$
\begin{tabular}{c|c}
	12 & 0\\
	6 & \\
\end{tabular}
$\rightarrow$
\begin{tabular}{c|c}
	12 & 0\\
	6 & 0\\
	3 & \\
\end{tabular}
$\rightarrow$
\begin{tabular}{c|c}
	12 & 0\\
	6 & 0\\
	3 & 1\\
	1 & \\
\end{tabular}
$\rightarrow$
\begin{tabular}{c|c}
	12 & 0\\
	6 & 0\\
	3 & 1\\
	1 & 1\\
	0 &
\end{tabular}

$12_{10} = 1100_2$
\end{frame}


\begin{frame}{Esempi conversione da base 10 a base 2}
Per controllare se si sono fatti errori si può controllare che i resti valgano 1 solo in corrispondenza di numeri dispari.
Inoltre, convertendo un numero $N$, il numero di cifre deve essere inferiore a $\ceil{log_2N}$. Equivalentemente, detto $C$ il numero di cifre utilizzate, si dovrà avere $2^M \geq N$ e $2^{M-1} < N$

\begin{tabular}{c|c}
	27 & 1\\
	13 & 1\\
	6 & 0\\
	3 & 1\\
	1 & 1\\
	0 &
\end{tabular}
\hskip 1cm
\begin{tabular}{c|c}
	35 & 1\\
	17 & 1\\
	8 & 0\\
	4 & 0\\
	2 & 0\\
	1 & 1\\
	0 & \\
\end{tabular}
\hskip 1cm
\begin{tabular}{c|c}
	42 & 0\\
	21 & 1\\
	10 & 0\\
	5 & 1\\
	2 & 0\\
	1 & 1\\
	0 & \\
\end{tabular}
\hskip 1cm
\begin{tabular}{c|c}
	73 & 1\\
	36 & 0\\
	18 & 0\\
	9 & 1\\
	4 & 0\\
	2 & 0\\
	1 & 1\\
	0 & \\
\end{tabular}
\end{frame}

\begin{frame}{Somme e sottrazioni}
Ragionando su singoli bit, le operazioni di somma e sottrazione sono definite nel seguente modo:

\begin{table}[h]
	\begin{tabular}{c|c}
		0+0 & =0 \\
		0+1 & =1 \\
		1+0 & =1 \\
		1+1 & =0 con riporto di 1
	\end{tabular}
\end{table}
\begin{table}[h]
	\begin{tabular}{c|c}
		0-0 & =0 \\
		1-0 & =1 \\
		1-1 & =0 \\
		0-1 & =1 con prestito di 1
	\end{tabular}
\end{table}

(si, sono giuste)
\end{frame}

\begin{frame}{Esercizi di somma}
Le somme si eseguono come in base 10. Se alla fine dell'operazione ``avanza'' un riporto, \alert{non possiamo} inserire un nuovo bit per tenerne conto, ma dobbiamo segnalare che si è verificato un \emph{overflow}, ovvero una condizione di errore.

\begin{tabular}{cc|c}
	primo operando   & 1001 & + \\
	secondo operando & 0101 & = \\
	riporti          &      &  \\
	\hline
	risultato        &   & \\
\end{tabular}
$\rightarrow$
\begin{tabular}{c|c}
	$1001$ & + \\
	$0101$ & = \\
	$\ \ 1\ $ &   \\
	\hline
	$\ \ \ 0$&  \\
\end{tabular}
$\rightarrow$
\begin{tabular}{c|c}
	$1001$ & + \\
	$0101$ & = \\
	$\ 01\ $&   \\
	\hline
	$\ \ 10$&  \\
\end{tabular}

$\rightarrow$
\begin{tabular}{c|c}
	$1001$ & + \\
	$0101$ & = \\
	$001\ $&   \\
	\hline
	$\ 110$&  \\
\end{tabular}
$\rightarrow$
\begin{tabular}{c|c}
	$1001$ & + \\
	$0101$ & = \\
	$001\ $&   \\
	\hline
	$1110$&  \\
\end{tabular}
\end{frame}

\begin{frame}{Esercizi di somma}
\begin{tabular}{c|c}
	$1101$ & + \\
	$0101$ & = \\
	\hskip-0.25cm\alert{1}$101\ $&   \\
	\hline
	$0010$& \alert{OVF}\\
\end{tabular}
\hskip 1.5cm
\begin{tabular}{c|c}
	$0110$ & + \\
	$1010$ & = \\
	\hskip-0.25cm\alert{1}$110\ $&   \\
	\hline
	$0000$& \alert{OVF}\\
\end{tabular}
\hskip 1.5cm
\begin{tabular}{c|c}
	$0110$ & + \\
	$0111$ & = \\
	\hskip-0.1cm$110\ $&   \\
	\hline
	$1101$& \\
\end{tabular}

\vskip 1cm
\begin{tabular}{c|c}
	$10011$ & + \\
	$01101$ & = \\
	\hskip-0.25cm\alert{1}$1111\ $&   \\
	\hline
	$00000$& \alert{OVF} \\
\end{tabular}
\hskip 1.5cm
\begin{tabular}{c|c}
	$010101$ & + \\
	$001010$ & = \\
	\hskip-0.1cm$00000\ $&   \\
	\hline
	$011111$& \\
\end{tabular}
\hskip 1.5cm
\begin{tabular}{c|c}
	$010101$ & + \\
	$001111$ & = \\
	\hskip-0.1cm$11111\ $&   \\
	\hline
	$100100$& \\
\end{tabular}
\end{frame}

\begin{frame}{Controllo esercizi di somma}
Convertiamo le somme in base 10 per trovare eventuali errori:

\begin{itemize}
\item $1+4+16 \ \ + \ \ 1+4 = 21+5 = 26 > 2^4 \rightarrow$ \alert{OVF}
\item $2+4 \ \ + \ \ 2+8 = 6+10 = 16 = 2^4 \rightarrow$ \alert{OVF} (con 4 bit rappresentiamo 16 valori, ma considerando lo zero 16 diventa il \emph{diciassettesimo})
\item $2+4 \ \ + \ \ 1+2+4 = 6+7 = 13 = 1+4+8$
\item $1+2+16 \ \ + \ \ 1+4+8 = 19+13 = 32 = 2^5 \rightarrow$ \alert{OVF}
\item $1+4+16 \ \ + \ \ 2+8 = 21+10 = 31 = 1+2+4+8+16$
\item $1+4+16 \ \ + \ \ 1+2+4+8 = 21+15 = 36 = 4+32$
\end{itemize}
\end{frame}
\end{document}